\section{Additional Results and Diagnostics}
\label{sec:appendix-results}

\subsection{Raw validity across candidate budgets}
\label{sec:appendix-validity}

Table~\ref{tab:raw-validity} reports raw candidate validity separately from
strict property success. Group-RL improves in-distribution validity by roughly
five points, but OOD validity remains close to 12\% for both policies. Nearly
every valid draw is unique, indicating that invalid SMILES syntax rather than
conventional mode collapse is the primary sample-efficiency bottleneck.

\begin{table}[h]
  \centering
  \caption{Raw candidate validity (\%) at matched generation prefixes.}
  \label{tab:raw-validity}
  \small
  \begin{tabular}{llrrr}
    \toprule
    Evaluation & Method & $k=1$ & $k=8$ & $k=20$ \\
    \midrule
    2p--7p & SFT & 29.20 & 28.85 & 29.04 \\
    2p--7p & Group-RL & \textbf{33.53} & \textbf{33.86} & \textbf{34.25} \\
    OOD & SFT & \textbf{12.50} & \textbf{11.92} & 11.51 \\
    OOD & Group-RL & 12.10 & 11.63 & \textbf{11.60} \\
    \bottomrule
  \end{tabular}
\end{table}

\subsection{Larger rollout groups are not sufficient}
\label{sec:appendix-g64}

Table~\ref{tab:g64-ablation} compares the main group-16 policy with a
hard-only continuation using 64 rollouts per program on a fixed 6p/7p subset.
The larger group improves validity but reduces Pass@20 from 50.78\% to 45.31\%,
driven by a large 7p regression. We therefore retain group 16 and report group
64 as a negative ablation rather than selecting it on raw validity.

\begin{table}[h]
  \centering
  \caption{Hard-program rollout-group ablation (\%). G64 continues training
  from the G16 checkpoint using only 6p/7p programs.}
  \label{tab:g64-ablation}
  \scriptsize
  \setlength{\tabcolsep}{3.5pt}
  \begin{tabular}{lrrrrrr}
    \toprule
    Method & $k$ & Raw strict & Pass@$k$ & 6p Pass@$k$ & 7p Pass@$k$ & Validity \\
    \midrule
    SFT & 8 & 3.47 & 21.88 & 21.88 & 21.88 & 28.27 \\
    Group-RL G16 & 8 & \textbf{4.74} & \textbf{28.12} & 31.25 & \textbf{25.00} & 32.91 \\
    Group-RL G64 & 8 & 4.44 & 25.78 & \textbf{32.81} & 18.75 & \textbf{36.38} \\
    \midrule
    SFT & 20 & 3.63 & 37.11 & 39.84 & 34.38 & 29.00 \\
    Group-RL G16 & 20 & 4.75 & \textbf{50.78} & \textbf{50.78} & \textbf{50.78} & 33.26 \\
    Group-RL G64 & 20 & \textbf{5.02} & 45.31 & \textbf{50.78} & 39.84 & \textbf{35.63} \\
    \bottomrule
  \end{tabular}
\end{table}

\subsection{MolEdit-Instruct comparison}
\label{sec:appendix-moledit}

Table~\ref{tab:moledit-horizontal} places MolProgram alongside seven published
MolEdit-Instruct baselines. External values are ten-task macro-averages
recomputed from Table~1 of \citet{zhuang2025moleditrl}. The comparison exposes
the boundary of direct decoding: MolProgram's first raw candidate is much less
valid and less source-preserving than every specialized editing model. A
property-aware 20-candidate selector improves relaxed accuracy but remains
assisted inference and does not repair the strict 0.65 threshold.

\begin{table}[h]
  \centering
  \caption{MolEdit-Instruct ten-task macro-average (\%).
  $\operatorname{Acc}_{\mathrm{all}}(t)$ requires both the requested property
  directions and source Tanimoto similarity of at least $t$.}
  \label{tab:moledit-horizontal}
  \scriptsize
  \setlength{\tabcolsep}{4pt}
  \begin{tabular}{llrrr}
    \toprule
    Method & Protocol & Validity & $\operatorname{Acc}_{\mathrm{all}}(.65)$ & $\operatorname{Acc}_{\mathrm{all}}(.15)$ \\
    \midrule
    BioT5 & reported & \textbf{100.0} & 0.0 & 29.2 \\
    MolGen & reported & \textbf{100.0} & 3.3 & 22.2 \\
    REINVENT4 & reported & 64.1 & 12.5 & 25.0 \\
    DrugAssist & reported & 97.9 & 34.9 & 41.1 \\
    GeLLM3O$_{\mathrm{Mistral}}$ & reported & 90.2 & 15.2 & 35.4 \\
    GeLLM3O$_{\mathrm{Llama}}$ & reported & 90.1 & 14.5 & 40.4 \\
    MolEditRL & reported & 96.6 & \textbf{45.0} & \textbf{72.7} \\
    \midrule
    MolProgram (Group-RL) & first raw candidate & 23.4 & 0.0 & 3.9 \\
    MolProgram (SFT) & assisted, $k=20$ & 96.4 & 0.0 & 17.9 \\
    MolProgram (Group-RL) & assisted, $k=20$ & 97.2 & 0.0 & 19.0 \\
    \bottomrule
  \end{tabular}
\end{table}

\subsection{MuMO source-fidelity diagnostic}
\label{sec:appendix-mumo}

Under the official 20-candidate MuMO oracle, MolProgram reaches 50.4\% overall
property SR, with 32.3\% IND and 68.8\% OOD SR. However, mean source
similarity among successful outputs is only 0.098, and no input passes the
additional 0.4 source-similarity diagnostic. The policy therefore often finds
a molecule with the desired properties by changing the molecular identity
rather than editing the input locally. MuMO SR is useful evidence for property
control, but it must not be presented as source-preserving editing success.

\begin{table}[h]
  \centering
  \caption{MolProgram on MuMOInstruct under the official property SR and an
  additional source-fidelity diagnostic.}
  \label{tab:mumo-fidelity}
  \small
  \begin{tabular}{lrrrrr}
    \toprule
    Method & Overall SR & IND SR & OOD SR & Sim(success) & Sim$\geq .4$ \\
    \midrule
    MolProgram (Group-RL) & 50.4 & 32.3 & 68.8 & 0.098 & 0.0 \\
    \bottomrule
  \end{tabular}
\end{table}

\subsection{Executable actions improve editing reachability}
\label{sec:appendix-actions}

A source-similarity reward applied to free-form direct SMILES Group-RL did not
produce any candidates above the required source-similarity threshold. In a
paired 200-example validation experiment, replacing free-form editing with
executable GraphEditDSL actions raises raw $k=1$ validity from 31.0\% to
100\% and strict $\operatorname{Acc}_{\mathrm{all}}(.65)$ from 0\% to 20.0\%,
without reducing held-out de novo strict success. At $k=8$, an explicitly
assisted finalizer reaches 43.5\% strict accuracy. These validation results
motivate source-consistent action policies, but are not part of the main
MolProgram claim.

\subsection{Frozen slots for ongoing repair experiments}
\label{sec:appendix-pending}

We retain empty cells for the preregistered repair runs rather than filling
them with selected-molecule validity or results from a different evaluation
subset.

\begin{table}[h]
  \centering
  \caption{Syntax-safe decoding repair. Pending rows remain blank until the
  paired P2 evaluation completes.}
  \label{tab:validity-repair}
  \small
  \begin{tabular}{llrrr}
    \toprule
    Evaluation & Decoder & Raw validity@1 & Raw strict@1 & Pass@8 \\
    \midrule
    2p--7p & legacy Group-RL & 33.53 & 8.93 & 41.72 \\
    2p--7p & syntax-safe Group-RL & -- & -- & -- \\
    OOD & legacy Group-RL & 12.10 & 3.20 & 23.60 \\
    OOD & syntax-safe Group-RL & -- & -- & -- \\
    \bottomrule
  \end{tabular}
\end{table}

\begin{table}[h]
  \centering
  \caption{Source-consistent executable editing. Pending rows use source-only
  consistency ranking; target molecules and output property oracles do not
  enter raw selection.}
  \label{tab:edit-repair}
  \scriptsize
  \setlength{\tabcolsep}{3pt}
  \begin{tabular}{lrrr}
    \toprule
    Edit policy & Validity@1 & $\operatorname{Acc}_{\mathrm{all}}(.65)$ & $\operatorname{Acc}_{\mathrm{all}}(.15)$ \\
    \midrule
    Protected action policy & 100.0 & 20.0 & 20.0 \\
    Source-consistent action policy & -- & -- & -- \\
    Strong source-consistent action policy & -- & -- & -- \\
    \bottomrule
  \end{tabular}
\end{table}
